\documentclass[10pt,a4paper]{amsart}
\usepackage[margin=19mm]{geometry}
\usepackage{amsmath,amssymb,mathtools}
\usepackage[scr=rsfs,cal=euler]{mathalfa}
\usepackage{xfrac}
\usepackage[unicode,hidelinks,bookmarks=false]{hyperref}
\newcommand{\ie}{i.e.\ }
\newcommand{\cf}{cf.\ }
\newcommand{\resp}{respectively\ }
\newcommand{\Subsection}{\subsection}
\providecommand{\nobreakdash}{\nobreak\hbox{-}}
\setlength{\emergencystretch}{2em}
\multlinegap0pt
\usepackage{enumerate}
\usepackage{amssymb}
\usepackage[shortlabels]{enumitem}
\usepackage{tikz}
\newtheorem{assumption}{Assumption}
\newtheorem{lemma}{Lemma}
\newtheorem{theorem}{Theorem}
\def\geq{\geqslant}
\def\leq{\leqslant}
\title{\LARGE Quantitative Error Estimates for Learning Macroscopic Mobilities from Microscopic Fluctuations}
\author{Nicolas Dirr, Zhengyan Wu,, Johannes Zimmer}
\date{Source: arXiv:2603.05089v1 (CC BY 4.0)}
\begin{document}
\maketitle

\noindent\emph{Source and licence.} \LARGE Quantitative Error Estimates for Learning Macroscopic Mobilities from Microscopic Fluctuations. Version arXiv:2603.05089v1 (CC BY 4.0), distributed by arXiv under the Creative Commons Attribution 4.0 International licence, http://creativecommons.org/licenses/by/4.0/, as recorded in the arXiv licence metadata for this exact version and shown on the abstract page https://arxiv.org/abs/2603.05089v1. Full licence notice: \texttt{licenses/arxiv-2603.05089-CC-BY-4.0.txt}.\par\smallskip

\noindent\textit{Editorial scope.} Counted unit: the article's theorem theorem (source0.tex lines 1597--1613). The article's complete original proof of it is delivered with it (source0.tex lines 1614--1825, one \texttt{proof} environment of 212 lines). No mathematics is written, repaired or re-derived: the statement and the proof are the article's own lines, in the article's own notation, and no retained line is substituted or re-flowed.  Retained uncounted as the source-local context the proof needs: lines 797--799; lines 807--816; lines 1540--1542; lines 1545--1547; lines 1570--1581.  Nothing was omitted from any retained span, so the package carries no omission record.  No retained line contains a citation, so the wrapper carries no bibliography and no bibliography entry.  No retained line contains a figure environment, an image-inclusion command or drawing code, so no image is delivered and the package declares no asset dependency.  Editorial formatting only, in addition: an AMS \texttt{amsart} preamble that restates the article's own declarations and macros used by the retained lines, namely the environments \texttt{assumption}, \texttt{lemma}, \texttt{theorem} on separate counters, together with the standard packages those lines need.  The retained environments print in this document as Assumption 1, Lemma 1, Theorem 1 in order of appearance, whereas the article prints the same items with the numbering of its own full document.  The complete native source of the article as distributed in the arXiv e-print for this exact version is bundled unchanged as \texttt{sources/195783/source0.tex}.
\begin{assumption}\label{Assump-initial-1}
We assume that the initial profile satisfies $\rho_0 \in C^{\infty}(\mathbb{T}^d)$ and is uniformly bounded away from the boundary values $0$ and $1$, i.e., there exists $c \in (0,1/2)$ such that $c \leq \rho_0 \leq 1 - c$ almost everywhere.
\end{assumption}

\begin{assumption}\label{Assump-sigman}
Suppose $\sigma(\cdot) \in C^{1}_{c}((0,1))$ and satisfies the following properties:
\begin{enumerate}
    \item[(1)] $\sigma \in C([0,1)) \cap C^{\infty}((0,1))$ with $\sigma(0)=\sigma(1)=0$ and $\sigma' \in C^{\infty}_c([0,1))$;
    \item[(2)] There exists a constant $c \in (0,\infty)$ such that for all $\zeta \in [0,1]$,
    \begin{align}
        |\sigma(\zeta)| \leq c, \quad \text{and} \quad |\sigma(\zeta)\sigma'(\zeta)| \leq c.
    \end{align}
\end{enumerate}
\end{assumption}

\begin{align}\label{SPDE-1-ito-sec-5}
d\rho_{\varepsilon} = \frac{1}{2}\Delta \rho_{\varepsilon}\,dt - \varepsilon^{1/2} \nabla\cdot\left(\sigma_{n(\varepsilon)}(\rho_{\varepsilon})\, dW_{\delta(\varepsilon)}\right) + \frac{\varepsilon}{2} \nabla\cdot\left(\sigma_{n(\varepsilon)}'(\rho_{\varepsilon})^2\, \nabla \rho_{\varepsilon} \right) F_{1,\delta(\varepsilon)},\quad \rho_{\varepsilon}(0)=\rho_0. 
\end{align}

\begin{equation}\label{PDE}
d\bar{\rho} = \frac{1}{2}\Delta\bar{\rho}\, dt, \quad \bar{\rho}(0) = \rho_0.
\end{equation}

\begin{lemma}\label{lem-LLN-regular}
Assume that the initial condition $\rho_0$ and the regularization sequence $(\sigma_n(\cdot))_{n \geq 1}$ satisfy Assumptions~\ref{Assump-initial-1} and~\ref{Assump-sigman}, respectively. For each $\varepsilon > 0$, let $\rho_{\varepsilon}$ denote the mild solution to \eqref{SPDE-1-ito-sec-5} with initial condition $\rho_0$, and let $\bar{\rho}$ denote the mild solution to \eqref{PDE} with the same initial data. Then, under the scaling condition
\begin{align*}
\lim_{\varepsilon \to 0} \varepsilon \delta(\varepsilon)^{-d-2} = 0,
\end{align*}
it follows that
\begin{align*}
\sup_{t\in[0,T]}\mathbb{E} \left\| \rho_{\varepsilon}(t) - \bar{\rho}(t) \right\|_{L^2(\mathbb{T}^d)}^2 
\leq \varepsilon \delta(\varepsilon)^{-d-2} 
\longrightarrow 0 \quad \text{as } \varepsilon \to 0.
\end{align*}
\end{lemma}

\begin{theorem}
Let $\sigma(\zeta)=\sqrt{\zeta(1-\zeta)}$. Assume that the initial condition $\rho_0$ satisfies Assumption \ref{Assump-initial-1} and the regularization sequence $(\sigma_n(\cdot))_{n \geq 1}$ satisfies Assumption \ref{Assump-sigman} with $\sigma_n\rightarrow\sigma$ in $C^1_{loc}((0,1))$. Let $c$ denote the lower bound of the initial data, as specified in Assumption \ref{Assump-initial-1}. Let $\phi \in C^{\infty}(\mathbb{T}^d)$, and define $\rho_{\varepsilon}^{\phi} := \langle \rho_{\varepsilon}, \phi \rangle$. The corresponding fluctuation field is given by
\begin{align*}
\bar{\rho}_{\varepsilon}^{1,\phi} := \varepsilon^{-1/2}(\rho_{\varepsilon}^{\phi} - \bar{\rho}^{\phi}).
\end{align*}
Then, for every $t > h > 0$, the following estimate holds:
\begin{align*}
\left| \frac{1}{h} \mathbb{E} \left( \bar{\rho}_{\varepsilon}^{1,\phi}(t+h) - \bar{\rho}_{\varepsilon}^{1,\phi}(t) \right)^2 
- \left\langle \nabla \phi, \bar{\rho}(t)(1 - \bar{\rho}(t)) \nabla \phi \right\rangle \right|
\leq C(\phi, \rho_0, t) \cdot \mathrm{Error}(h, \varepsilon, \delta(\varepsilon), n(\varepsilon)),
\end{align*}
where the error term is given by
\begin{align*}
\mathrm{Error}(h, \varepsilon, \delta(\varepsilon), n(\varepsilon)) =\ & h 
+ \varepsilon \delta(\varepsilon)^{-2d} \|\sigma_{n(\varepsilon)}'\|_{L^{\infty}([0,1])}^4 h + \varepsilon \delta(\varepsilon)^{-d-2}. 
\end{align*}
\end{theorem}

\begin{proof}
We expand the difference $\bar{\rho}^{1,\phi}_{\varepsilon}(t+h) - \bar{\rho}^{1,\phi}_{\varepsilon}(t)$ using Duhamel's formula. The integral is naturally split into two parts: one over the interval $[0, t]$ and the other over $[t, t+h]$. A straightforward computation yields the following expression:
\begin{align*}
&\bar{\rho}^{1,\phi}_{\varepsilon}(t+h) - \bar{\rho}^{1,\phi}_{\varepsilon}(t)\\ 
=& - \int_0^{t+h} \Big\langle \nabla\phi, S(t+h - s) \sigma_{n(\varepsilon)}(\rho_{\varepsilon}) \, dW_{\delta(\varepsilon)}(s) \Big\rangle 
+ \int_0^t \Big\langle \nabla\phi, S(t - s) \sigma_{n(\varepsilon)}(\rho_{\varepsilon}) \, dW_{\delta(\varepsilon)}(s) \Big\rangle \\
& + \frac{\varepsilon^{1/2}}{2} F_{1,\delta(\varepsilon)} \int_0^{t+h} \left\langle \nabla\phi, S(t+h - s) \left( \sigma_{n(\varepsilon)}'(\rho_{\varepsilon})^2 \nabla \rho_{\varepsilon} \right) \right\rangle ds \\
& - \frac{\varepsilon^{1/2}}{2} F_{1,\delta(\varepsilon)} \int_0^t \left\langle \nabla\phi, S(t - s) \left( \sigma_{n(\varepsilon)}'(\rho_{\varepsilon})^2 \nabla \rho_{\varepsilon} \right) \right\rangle ds \\
=& - \int_0^t \Big\langle S(h)\nabla\phi - \nabla\phi, S(t - s) \sigma_{n(\varepsilon)}(\rho_{\varepsilon}) \, dW_{\delta(\varepsilon)}(s) \Big\rangle 
- \int_t^{t+h} \Big\langle \nabla\phi, S(t+h - s) \sigma_{n(\varepsilon)}(\rho_{\varepsilon}) \, dW_{\delta(\varepsilon)}(s) \Big\rangle \\
& + \frac{\varepsilon^{1/2}}{2} F_{1,\delta(\varepsilon)} \int_0^{t+h} \left\langle \nabla\phi, S(t+h - s) \left( \sigma_{n(\varepsilon)}'(\rho_{\varepsilon})^2 \nabla \rho_{\varepsilon} \right) \right\rangle ds \\
& - \frac{\varepsilon^{1/2}}{2} F_{1,\delta(\varepsilon)} \int_0^t \left\langle \nabla\phi, S(t - s) \left( \sigma_{n(\varepsilon)}'(\rho_{\varepsilon})^2 \nabla \rho_{\varepsilon} \right) \right\rangle ds.
\end{align*}
Taking the second moment, expanding the square, and applying It\^o's product rule, we observe that all cross-covariance terms vanish. Using It\^o's isometry along with the covariance structure of the Brownian motions, we obtain
\begin{align*}
\mathbb{E} \left( \bar{\rho}^{1,\phi}_{\varepsilon}(t+h) - \bar{\rho}^{1,\phi}_{\varepsilon}(t) \right)^2 
=&\ \mathbb{E} \int_0^t \| (S(h)\nabla\phi - \nabla\phi)\, S(t - s) \sigma_{n(\varepsilon)}(\rho_{\varepsilon}) \|_{L^2(\mathbb{T}^d)}^2\, ds \\
&+ \mathbb{E} \int_t^{t+h} \| \nabla\phi\, S(t+h - s)\, \sigma_{n(\varepsilon)}(\rho_{\varepsilon}) \|_{L^2(\mathbb{T}^d)}^2\, ds \\
&+ \frac{\varepsilon}{4} F_{1,\delta(\varepsilon)}^2\, \mathbb{E} \left( \int_0^{t+h} \left\langle \nabla\phi, S(t+h - s) \left( \sigma_{n(\varepsilon)}'(\rho_{\varepsilon})^2 \nabla \rho_{\varepsilon} \right) \right\rangle ds \right. \\
&\qquad \left. - \int_0^t \left\langle \nabla\phi, S(t - s) \left( \sigma_{n(\varepsilon)}'(\rho_{\varepsilon})^2 \nabla \rho_{\varepsilon} \right) \right\rangle ds \right)^2 \\
=:&\ I_1(h,\varepsilon) + I_2(h,\varepsilon) + I_3(h,\varepsilon).
\end{align*}

We now estimate each of the terms $I_1(h,\varepsilon)$, $I_2(h,\varepsilon)$, and $I_3(h,\varepsilon)$ separately. For the first term $I_1(h,\varepsilon)$, applying the convolutional Young's inequality yields
\begin{align*}
I_1(h,\varepsilon) 
&\leq \mathbb{E} \int_0^t \| S(h)\nabla\phi - \nabla\phi \|_{L^\infty(\mathbb{T}^d)}^2 
\cdot \| S(t - s)\sigma_{n(\varepsilon)}(\rho_{\varepsilon}) \|_{L^{2}(\mathbb{T}^d)}^2\, ds \\
&\leq C(\rho_0)\int_0^t \Big(\int^h_0\| S(r)\Delta\nabla\phi \|_{L^\infty(\mathbb{T}^d)}dr\Big)^2 
\cdot \| p(t - s) \|_{L^1(\mathbb{T}^d)}^2\, ds \\
&\leq C(\phi,\rho_0)\, h^2 t,
\end{align*}
where $p(t - s)$ denotes the heat kernel on $\mathbb{T}^d$ and the final inequality follows from the smoothing property of the heat semigroup.


To estimate $I_2(h,\varepsilon)$, we decompose the integrand by introducing an intermediate term involving the pointwise evaluation at time $t$. Specifically, we write:
\begin{align*}
I_2(h,\varepsilon) 
=&\ \mathbb{E} \int_t^{t+h} \left( \| \nabla\phi\, S(t+h-s)\, \sigma_{n(\varepsilon)}(\rho_{\varepsilon}(s)) \|_{L^2(\mathbb{T}^d)}^2 
- \| \nabla\phi\, \sigma_{n(\varepsilon)}(\rho_{\varepsilon}(t)) \|_{L^2(\mathbb{T}^d)}^2 \right) ds \\
&+ \mathbb{E} \int_t^{t+h} \left( \| \nabla\phi\, \sigma_{n(\varepsilon)}(\rho_{\varepsilon}(t)) \|_{L^2(\mathbb{T}^d)}^2 
- \| \nabla\phi\, \sqrt{\bar{\rho}(t)(1 - \bar{\rho}(t))} \|_{L^2(\mathbb{T}^d)}^2 \right) ds \\
&+ h \| \nabla\phi\, \sqrt{\bar{\rho}(t)(1 - \bar{\rho}(t))} \|_{L^2(\mathbb{T}^d)}^2 \\
=:&\ I_{2,1}(h,\varepsilon) + I_{2,2}(h,\varepsilon) + h \| \nabla\phi\, \sqrt{\bar{\rho}(t)(1 - \bar{\rho}(t))} \|_{L^2(\mathbb{T}^d)}^2.
\end{align*}

We first consider $I_{2,1}(h,\varepsilon)$. Using the non-negativity of $\rho_{\varepsilon}$ and the preservation of mass, and expanding the square via the polarization identity, we obtain:
\begin{align*}
I_{2,1}(h,\varepsilon) 
=&\ \mathbb{E} \int_t^{t+h} \int_{\mathbb{T}^d} \left( 
|\nabla\phi\, S(t+h - s)\, \sigma_{n(\varepsilon)}(\rho_{\varepsilon}(s))|^2 
- |\nabla\phi\, \sigma_{n(\varepsilon)}(\rho_{\varepsilon}(t))|^2 
\right) dx\, ds \\
=&\ \mathbb{E} \int_t^{t+h} \int_{\mathbb{T}^d} \left( 
\nabla\phi\, S(t+h - s)\, \sigma_{n(\varepsilon)}(\rho_{\varepsilon}(s)) 
+ \nabla\phi\, \sigma_{n(\varepsilon)}(\rho_{\varepsilon}(t)) 
\right) \cdot \\
&\hspace{3cm} \left( 
\nabla\phi\, S(t+h - s)\, \sigma_{n(\varepsilon)}(\rho_{\varepsilon}(s)) 
- \nabla\phi\, \sigma_{n(\varepsilon)}(\rho_{\varepsilon}(t)) 
\right) dx\, ds \\
\leq&\ C(\phi) \mathbb{E} \int_t^{t+h} \| S(t+h - s)\, \sigma_{n(\varepsilon)}(\rho_{\varepsilon}(s)) 
- \sigma_{n(\varepsilon)}(\rho_{\varepsilon}(t)) \|_{L^2(\mathbb{T}^d)} ds.
\end{align*}

We then split the difference inside the $L^1$-norm into two terms:
\begin{align*}
I_{2,1}(h,\varepsilon) 
\leq&\ C(\phi) \mathbb{E} \int_t^{t+h} \| S(t+h - s)\, \sigma_{n(\varepsilon)}(\rho_{\varepsilon}(s)) 
- S(t+h - s)\, \sigma_{n(\varepsilon)}(\rho_{\varepsilon}(t)) \|_{L^2(\mathbb{T}^d)} ds \\
&+ C(\phi) \mathbb{E} \int_t^{t+h} \| S(t+h - s)\, \sigma_{n(\varepsilon)}(\rho_{\varepsilon}(t)) 
- \sigma_{n(\varepsilon)}(\rho_{\varepsilon}(t)) \|_{L^2(\mathbb{T}^d)} ds \\
=:&\ I_{2,1,1}(h,\varepsilon) + I_{2,1,2}(h,\varepsilon).
\end{align*}
We now estimate the term $I_{2,1,1}(h,\varepsilon)$. By applying convolutional Young's inequality to the heat semigroup $S(t)$, we obtain
\begin{align*}
I_{2,1,1}(h,\varepsilon) 
&\leq C(\phi) \mathbb{E} \int_t^{t+h} \| \sigma_{n(\varepsilon)}(\rho_{\varepsilon}(s)) - \sigma_{n(\varepsilon)}(\rho_{\varepsilon}(t)) \|_{L^2(\mathbb{T}^d)} ds.
\end{align*}
Using the Lipschitz continuity of $\sigma_{n(\varepsilon)}$ and the triangle inequality, we decompose the integrand into microscopic and macroscopic fluctuations:
\begin{align*}
I_{2,1,1}(h,\varepsilon) 
\leq\ & C(\phi) \| \sigma_{n(\varepsilon)}' \|_{L^{\infty}([0,1])} \mathbb{E} \int_t^{t+h} \| \rho_{\varepsilon}(s) - \bar{\rho}(s) \|_{L^2(\mathbb{T}^d)} ds \\
& + C(\phi) \| \sigma_{n(\varepsilon)}' \|_{L^{\infty}([0,1])} \mathbb{E} \| \rho_{\varepsilon}(t) - \bar{\rho}(t) \|_{L^2(\mathbb{T}^d)}\, h \\
& + C(\phi) \| \sigma_{n(\varepsilon)}' \|_{L^{\infty}([c,1-c])} \mathbb{E} \int_t^{t+h} \| \bar{\rho}(s) - \bar{\rho}(t) \|_{L^2(\mathbb{T}^d)} ds \\
=:&\ I_{2,1,1,1}(h,\varepsilon) + I_{2,1,1,2}(h,\varepsilon) + I_{2,1,1,3}(h,\varepsilon).
\end{align*}

To bound $I_{2,1,1,1}(h,\varepsilon)$, we apply Lemma \ref{lem-LLN-regular} to see that 
\begin{align*}
I_{2,1,1,1}(h,\varepsilon) 
&\leq C(\phi) \| \sigma_{n(\varepsilon)}' \|_{L^{\infty}([0,1])} 
\left( \varepsilon \delta(\varepsilon)^{-d-2} \right) h.
\end{align*}

For the second term, $I_{2,1,1,2}(h,\varepsilon)$, we use the same type of fluctuation estimate at a fixed time $t$:
\begin{align*}
I_{2,1,1,2}(h,\varepsilon) 
&\leq C(\phi) \left( \varepsilon \delta(\varepsilon)^{-d-2} \right) h.
\end{align*}

Finally, for the macroscopic contribution $I_{2,1,1,3}(h,\varepsilon)$, we use the temporal regularity of the macroscopic density $\bar{\rho}$, which is governed by the heat equation. In particular, the smoothing effect of the Laplacian implies that $\bar{\rho} \in W^{1,\infty}([0,T]; L^1(\mathbb{T}^d))$. Therefore, we obtain
\begin{align*}
I_{2,1,1,3}(h,\varepsilon) 
&\leq C(\phi) \| \sigma_{n(\varepsilon)}' \|_{L^{\infty}([c,1-c])} 
\| \bar{\rho} \|_{W^{1,\infty}([0,T]; L^2(\mathbb{T}^d))} h^2 \\
&\leq C(\phi, c) \| \bar{\rho} \|_{L^{\infty}([0,T]; H^2(\mathbb{T}^d))} h^2 
\leq C(\phi, c, \rho_0) h^2.
\end{align*}




We next estimate the term $I_{2,1,2}(h,\varepsilon)$. We have the following decomposition: 
\begin{align*}
I_{2,1,2}(h,\varepsilon)
\leq\ & C(\phi) \mathbb{E} \int_t^{t+h} \| S(t+h-s)\sigma_{n(\varepsilon)}(\rho_{\varepsilon}(t)) - S(t+h-s)\sigma_{n(\varepsilon)}(\bar{\rho}(t)) \|_{L^2(\mathbb{T}^d)} ds \\
&+ C(\phi) \mathbb{E} \int_t^{t+h} \| S(t+h-s)\sigma_{n(\varepsilon)}(\bar{\rho}(t)) - \sigma_{n(\varepsilon)}(\bar{\rho}(t)) \|_{L^2(\mathbb{T}^d)} ds\\
&+C(\phi)\mathbb{E}\int^{t+h}_t\|\sigma_{n(\varepsilon)}(\bar{\rho}(t))-\sigma_{n(\varepsilon)}(\rho_{\varepsilon}(t))\|_{L^2(\mathbb{T}^d)}ds. 
\end{align*}
Each term is estimated as follows, the first term and the third term can be controlled via the law of large numbers and Lipschitz continuity of $ \sigma_{n(\varepsilon)} $. The second term is controlled by the strong continuity of the semigroup $ S(t) $. Putting these together, applying H\"older's inequality and Lemma \ref{lem-LLN-regular}, we obtain
\begin{align*}
I_{2,1,2}(h,\varepsilon)\leq&C(\phi)\mathbb{E}\int^{t+h}_t\| \sigma_{n(\varepsilon)}(\bar{\rho}(t)) - \sigma_{n(\varepsilon)}(\rho_{\varepsilon}(t)) \|_{L^1(\mathbb{T}^d)} ds+C(\phi) h^2\\
\leq&C(\phi) \|\sigma_{n(\varepsilon)}'\|_{L^{\infty}([0,1])}(\sup_{t\in[0,T]}\mathbb{E}\|\rho_{\varepsilon}(t)-\bar{\rho}(t)\|_{L^2(\mathbb{T}^d)})h+C(\phi) h^2\\
\leq& C(\phi) \|\sigma_{n(\varepsilon)}'\|_{L^{\infty}([0,1])} 
\left( \varepsilon \delta(\varepsilon)^{-d-2} \right) h 
+ C(\phi) h^2.
\end{align*}

\medskip

We now consider $ I_{2,2}(h,\varepsilon) $. By the triangle inequality, we write
\begin{align*}
I_{2,2}(h,\varepsilon)
&\leq C(\phi) \mathbb{E} \int_t^{t+h} \int_{\mathbb{T}^d} \left| \sigma_{n(\varepsilon)}(\rho_{\varepsilon}(s))^2 - \bar{\rho}(t)(1 - \bar{\rho}(t)) \right| dx ds \\
&\leq C(\phi) \mathbb{E} \int_t^{t+h} \int_{\mathbb{T}^d} \left| \sigma_{n(\varepsilon)}(\rho_{\varepsilon}(s))^2 - \sigma_{n(\varepsilon)}(\bar{\rho}(s))^2 \right| dx ds \\
&\quad + \int_t^{t+h} \int_{\mathbb{T}^d} \left| \sigma_{n(\varepsilon)}(\bar{\rho}(s))^2 - \sigma_{n(\varepsilon)}(\bar{\rho}(t))^2 \right| dx ds \\
&\quad + h \int_{\mathbb{T}^d} \left| \sigma_{n(\varepsilon)}(\bar{\rho}(t))^2 - \bar{\rho}(t)(1 - \bar{\rho}(t)) \right| dx \\
&=: I_{2,2,1}(h,\varepsilon) + I_{2,2,2}(h,\varepsilon) + I_{2,2,3}(h,\varepsilon).
\end{align*}

For $ I_{2,2,1}(h,\varepsilon) $, we apply the Lipschitz bound on $ \sigma_{n(\varepsilon)}^2 $ and use the fluctuation estimate:
\begin{align*}
I_{2,2,1}(h,\varepsilon)
&\leq C(\phi) \| \sigma_{n(\varepsilon)} \sigma_{n(\varepsilon)}' \|_{L^{\infty}([0,1])} 
\mathbb{E} \int_t^{t+h} \| \rho_{\varepsilon}(s) - \bar{\rho}(s) \|_{L^1(\mathbb{T}^d)} ds \\
&\leq C(\phi) \| \sigma_{n(\varepsilon)} \sigma_{n(\varepsilon)}' \|_{L^{\infty}([0,1])} 
\left( \sup_{s \in [0,T]} \mathbb{E} \| \rho_{\varepsilon}(s) - \bar{\rho}(s) \|_{L^2(\mathbb{T}^d)}^2 \right)^{1/2} h \\
&\leq C(\phi) \left( \varepsilon \delta(\varepsilon)^{-d-2} \right) h.
\end{align*}

For $ I_{2,2,2}(h,\varepsilon) $, we exploit the temporal regularity of the macroscopic profile $ \bar{\rho} $:
\begin{align*}
I_{2,2,2}(h,\varepsilon)
&\leq C(\phi) \| \sigma_{n(\varepsilon)}\sigma_{n(\varepsilon)}' \|_{L^{\infty}([c,1-c])} 
\| \bar{\rho} \|_{W^{1,\infty}([0,T]; L^1(\mathbb{T}^d))} h^2 \\
&\leq C(\phi, c) \| \bar{\rho} \|_{L^{\infty}([0,T]; W^{2,1}(\mathbb{T}^d))} h^2 
\leq C(\phi, c, \rho_0) h^2.
\end{align*}

Lastly, for $ I_{2,2,3}(h,\varepsilon) $, the convergence of the approximation sequence $ \sigma_{n(\varepsilon)}^2 \to \zeta(1 - \zeta) $ implies
\begin{align*}
I_{2,2,3}(h,\varepsilon)
&\leq C(\phi) h \sup_{\zeta \in [c, 1 - c]} \left| \sigma_{n(\varepsilon)}(\zeta)^2 - \zeta(1 - \zeta) \right|.
\end{align*}
We observe that, the mollified noise coefficient $ \sigma_{n(\varepsilon)}(\cdot) $ is chosen such that, for sufficiently small $\varepsilon>0$, 
$$
\sigma_{n(\varepsilon)}(\zeta) = \sqrt{\zeta(1 - \zeta)} \quad \text{for all } \zeta \in [c, 1 - c],
$$
then for sufficiently small $ \varepsilon $, the approximation error satisfies
$$
\sup_{\zeta \in [c, 1 - c]} \left| \sigma_{n(\varepsilon)}(\zeta)^2 - \zeta(1 - \zeta) \right| = 0.
$$






Finally, we estimate the term $ I_3(h,\varepsilon) $. By applying integration by parts, we rewrite the expression as
\begin{align*}
I_3(h,\varepsilon)
= \frac{\varepsilon}{4} F_{1,\delta(\varepsilon)}^2 
\mathbb{E} \left( 
\int_0^{t+h} \langle \Delta \phi, S(t+h - s) \Sigma_{n(\varepsilon)}(\rho_{\varepsilon}(s)) \rangle ds 
- \int_0^t \langle \Delta \phi, S(t - s) \Sigma_{n(\varepsilon)}(\rho_{\varepsilon}(s)) \rangle ds 
\right)^2,
\end{align*}
where the function $ \Sigma_{n(\varepsilon)} \colon [0,1] \to \mathbb{R} $ is defined by
$$
\Sigma_{n(\varepsilon)}(\zeta) := \int_0^{\zeta} \sigma_{n(\varepsilon)}'(\zeta')^2 d\zeta'.
$$

Applying the triangle inequality and splitting the difference of integrals, we estimate:
\begin{align*}
I_3(h,\varepsilon)
&\lesssim \varepsilon \delta(\varepsilon)^{-2d} 
\mathbb{E} \left( \int_0^t \langle \Delta \phi, \left( S(t+h - s) - S(t - s) \right) \Sigma_{n(\varepsilon)}(\rho_{\varepsilon}(s)) \rangle ds \right)^2 \\
&\quad + \varepsilon \delta(\varepsilon)^{-2d} 
\mathbb{E} \left( \int_t^{t+h} \langle \Delta \phi, S(t+h - s) \Sigma_{n(\varepsilon)}(\rho_{\varepsilon}(s)) \rangle ds \right)^2.
\end{align*}

Since $ \Sigma_{n(\varepsilon)} $ is uniformly Lipschitz, and the heat semigroup is strongly continuous, both terms are controlled similarly. Using boundedness of $ \Delta \phi $ and standard smoothing properties of $ S(t) $, we obtain:
\begin{align*}
I_3(h,\varepsilon) 
\lesssim C(\phi) \varepsilon \delta(\varepsilon)^{-2d} 
\| \sigma_{n(\varepsilon)}' \|_{L^{\infty}([0,1])}^4 h^2.
\end{align*}

Combining the bounds on $ I_1(h,\varepsilon) $, $ I_2(h,\varepsilon) $, and $ I_3(h,\varepsilon) $, we conclude the proof.

\end{proof}

\end{document}
